\begin{figure}[H]
\centering
\begin{tikzpicture}[x=4.6cm,y=3.0cm,>=Stealth,font=\small]
\begin{scope}
\draw[->] (0,0)--(1.1,0) node[right] {$\kappa$};
\draw[->] (0,0)--(0,1.08) node[above] {$F(\kappa)$};
\draw[very thick,linkblue] (0,0)--(.36,0)--(1,.72);
\draw[dashed,gray] (1,0)--(1,.72)--(0,.72);
\node[left] at (0,.72) {$v$};
\node[below] at (.36,0) {$\tau$};
\node[below] at (1,0) {$1$};
\node[above left] at (.8,.5) {slope $\beta=v/\ell$};
\node[below=16pt] at (.55,0) {Canonical deficit profile};
\end{scope}
\begin{scope}[xshift=7.0cm]
\draw[->] (0,0)--(1.1,0) node[right] {$\kappa$};
\draw[->] (0,0)--(0,1.08) node[above] {$\kappa-F(\kappa)/v$};
\draw[very thick,linkblue] (0,0)--(.36,.55)--(1,0);
\draw[dashed,gray] (.36,0)--(.36,.55)--(0,.55);
\node[left] at (0,.55) {$\tau$};
\node[below] at (.36,0) {$\tau$};
\node[below] at (1,0) {$1$};
\node[above] at (.75,.3) {$\tau(1-\kappa)/\ell$};
\node[below=16pt] at (.55,0) {Remaining orientation exponent};
\end{scope}
\end{tikzpicture}
\caption{The canonical profile in Proposition~\ref{prop:canonical}.
After the flat interval, its slope is constant. For $\kappa\ge\tau$, the associated quantity
$\kappa-F(\kappa)/v$ decreases linearly from $\tau$ to zero;
at the stationary prefix $\kappa=1-\ell c$ it equals $\tau c$.
This identity gives the angular scale $\Theta_c=N^{-\tau c}$.
The vertical scales in the two panels are independent.}
\label{fig:canonical-profile}
\end{figure}
